\section{Comparability Audit Across All 58 Core Families}\label{sec:defenses}
\subsection{Unified 58-core disposition and 39-record ledger}
The corpus partitions as $82=58$ core $+24$ adjacent families. Within the core, $58=20$ mitigation $+6$ contextual $+32$ map-only families. The ledger holds 39 records: 33 from the 20 mitigation families and one from each contextual family---SIGLEAK, AgentTell, FP-Agent, Behind Bars, LLMs Have Rhythm and ML privacy channels \cite{geng2026sigleak,shahriar2026agenttell,wang2026fpagent,gu2026behindbars,alhazbi2025rhythm,debenedetti2024privacy}---which supply fixed-guardrail, access-control, interference, VPN or filter-origin contrasts rather than an audited mitigation claim. The 32 map-only families comprise 22 without a scoped mitigation evaluation, five diagnostics, four unbound quantitative assertions and SplitDecoy's boundary case. Table~\ref{tab:defense-types} groups mechanisms; its eight rows are mechanism groups, not sampled studies.
\begin{table*}[t]\centering\footnotesize
\caption{Seven privacy and attribution intervention categories organized by receiving view and supported deployment judgment. Type I protects confidentiality; Type II addresses provenance or identity attribution.}\label{tab:defense-types}
\begin{tabularx}{\textwidth}{@{}p{.17\textwidth}p{.29\textwidth}Y@{}}\toprule
Intervention \& polarity & Evaluated view, target and source & Established result and applicable boundary\\\midrule
Padding / batching (Type~I: confidentiality) & Network packet lengths and groups; topic inference. NetEcho and Whisper Leak \cite{netecho2025,whisperleak2025}. & Collected padding traces retain inference signal; packet-group batching is evaluated by replay. A protected upstream length view does not cover downstream re-emission or residual grouping. Live incremental cost remains \NRR{}.\\
Constant-rate / transport cover (Type~I: confidentiality) & Network workflow fingerprint; Dangol's OHTTP testbed \cite{dangol2026workflow}. & The co-located loopback setting yields 0.30 recovery at about $13\times$ bytes and $11\times$ latency. Transformed-capture results are separate; neither result supplies objective agent-task utility.\\
Concurrent query decoys (Type~I: confidentiality) & Destination/domain observer; private prompt and trait inference. WRA \cite{jeong2026wra}. & Real/decoy research prompts execute in parallel as read-only queries whose decoy answers are discarded locally; interleaved-domain evaluation reduces one-decoy $E_{\rm func}$ from 0.770 to 0.610. Terminal decoy utility/cost are \NRR{}.\\
Cache isolation / keying (Type~I: confidentiality) & Co-tenant prefix timing; SafeKV, KVGov and CachePrune \cite{chu2025safekv,addagada2026kvgov,wu2026cacheprune}. & SafeKV reports post-trigger AUC near 0.5 and separate TTFT/throughput comparisons. KVGov's HMAC median $1.3\,\mu$s prices key computation, excluding lost reuse. CachePrune jointly binds marked-token protection, QA/summarization utility and serving performance.\\
Cache timing equalization (Type~I: confidentiality) & Co-tenant cache-fusion timing; SpliceDefense \cite{spliceleak2026}. & Constant-Time Boundary Fusion (CTBF) suppresses the tested signal over 15 tasks/75 steps. Quantized Chunk Padding (QCP) adds mean 7.5 tokens and $<5$ms/$<1\%$ prefill time; this does not price combined deployment.\\
Timestamp perturbation (Type~II: attribution) & Recorded website action metadata; model/agent attribution. UI Traces \cite{uitraces2026}. & Fitting on delayed traces ($D_2$) recovers substantial attribution performance under offline timestamp edits that retain action types/order. Live execution delay and task outcomes are not measured by the trace transformation.\\
Oblivious / unshared memory (Type~I: confidentiality) & Host or local resource observations; TDXRay, Gao et al.\ and OTRO \cite{hornetz2026tdxray,gao2025iknow,lee2026otro}. & TDXRay measures oblivious-map lookup cost; Gao measures loading/memory cost after disabling zero-copy. OTRO reduces conditional access-count information (3.94 to 0.04 bits) with separate unpadded TTFT (4.59\%) and epoch-padding ($2.8$ms) bounds.\\\bottomrule
\end{tabularx}
\par\smallskip\raggedright\footnotesize
Authorized operator diagnostics (Type~IV: ClawGuard, cycle detection, LumiMAS, LLM-serving RCA and MIRAGE) are baseline telemetry rather than privacy attacks or evasion defenses, and are cataloged separately.
\end{table*}

\paragraph{Whole-ledger synthesis}
Each record has one primary measurement mode: 20 executed collections/prototypes, six trace transformations, five simulations/derived estimates, six component/calibration measurements and two unresolved placements (Table~\ref{tab:r9-record-index}). Eleven of the 20 mitigation families contribute more than one record, and nine of them span more than one mode; this is why a mechanism name cannot carry utility or cost between records. Executed collection establishes its stated execution scope, not automatically full-workflow utility.

\subsection{Within-family and within-cohort record comparisons}
The pairwise audit covers the 33 records in the ledger before the final search (27 mitigation, six contextual); the six records that search added contribute two further same-family pairs (\texttt{CACHEPRUNE-selective}/\texttt{-context} and \texttt{OTRO-unpadded}/\texttt{-epochpad}). Of the $\binom{33}{2}=528$ pairs, 426 cross unrelated mechanism cohorts and are out of scope by construction. The remaining 102 related pairs (66 traffic-envelope, 28 cache-service, three behavioral-humanization, three generator-fingerprint, one local-memory, one fixed-instruction) include the 13 same-family pairs. They show why paper-level comparison fails:
\begin{itemize}
\item \textbf{Same paper, different object.} Six of 13 same-family pairs (7 of 15 with the two later pairs) measure different primitives, targets or cost objects, e.g., Early Bird's $K$-prefix oracle simulation versus its GPTCache semantic prototype, or SpliceDefense's CTBF signal suppression versus QCP-only prefill cost.
\item \textbf{Same object, different population.} Of the seven aligned-outcome same-family pairs, four evaluate different collection or execution populations (e.g., Dangol's stored-capture ablation at 0.17 accuracy versus live loopback OHTTP at 0.30) and three lack a verified matched-test binding; none licenses subtraction as a paired deployment effect.
\item \textbf{Across papers.} Of the 102 related pairs, 16 (15.7\%) support a source-qualified descriptive comparison, 80 (78.4\%) have an established receiver, target or cost-object mismatch, and 6 (5.9\%) remain unresolved.
\end{itemize}

\subsection{Nine-item reuse audit across 20 mitigation families}
Table~\ref{tab:audit-rules} defines each item before its application to the 180 cells in Table~\ref{tab:family-audit}. Applicability follows the result's reuse claim: a fixed analytical attack makes adaptation inapplicable, while a component benchmark can support component utility. \NRR{} and \NRV{} retain distinct source states: the 11 unverified cells are five R4 and six R5 bindings, which are distinguished from an attack failure.
\begin{table}[t]\centering\footnotesize
\caption{Operational rules for the nine reuse bindings. Apply each item to a specified result record and deployment claim.}\label{tab:audit-rules}
\begin{tabularx}{\columnwidth}{@{}p{.30\columnwidth}Y@{}}\toprule
Item & Evidence required for a Bound (\textbf{B}) judgment\\\midrule
R1 Target & Hidden property, inference goal and authorized/protected interest are explicit.\\
R2 Receiver & Observer permission, released representation and auxiliary inputs identify the evaluated view.\\
R3 Intervention & Placement and execution mode bind the transform to stored data, wire emission or executed actions.\\
R4 Protected fit & Attacker fitting/calibration uses the protected condition; training and protected testing are identified.\\
R5 Matched adaptation & Clean- and protected-condition fits are evaluated on the same defended test population.\\
R6 Workload & Tasks, system/generator configuration and holdout specify the supported transfer axis.\\
R7 Inference & Metric, denominator and evaluation condition bind the attack outcome to the scoped mechanism.\\
R8 Utility & Task/component performance is measured at the same protected operating point, with its state boundary explicit.\\
R9 Cost & Resource units, baseline and measurement object price that operating point.\\\bottomrule
\end{tabularx}
\par\smallskip\raggedright\footnotesize \textbf{B}: bound; \textbf{P}: partial binding; \NRR{}: unreported in inspected scope; \NRA{}: source unavailable; \NRV{}: binding unverified; --: inapplicable. First determine applicability, then assign evidence status. Family \textbf{B} means at least one scoped record satisfies the item, not that every result does.
\end{table}

\begin{table}[t]\centering\footnotesize
\caption{Record-derived evidence availability for all 20 evaluated-mitigation families. R1--R9 are scoped reuse items, not overall paper scores.}\label{tab:family-audit}
\setlength{\tabcolsep}{2.2pt}
\begin{tabularx}{\columnwidth}{@{}Y*{9}{c}@{}}\toprule
Family & R1 & R2 & R3 & R4 & R5 & R6 & R7 & R8 & R9\\\midrule
UI Traces~\cite{uitraces2026} & \textbf{B} & \textbf{P} & \textbf{B} & \textbf{B} & \textbf{B} & \textbf{B} & \textbf{B} & U & --\\
WRA~\cite{jeong2026wra} & \textbf{B} & \textbf{P} & \textbf{B} & -- & -- & \textbf{P} & \textbf{B} & \textbf{P} & U\\
NetEcho~\cite{netecho2025} & \textbf{B} & \textbf{B} & \textbf{B} & \textbf{B} & ? & \textbf{B} & \textbf{B} & U & \textbf{P}\\
Early Bird~\cite{earlybird2025} & \textbf{B} & \textbf{B} & \textbf{B} & \textbf{B} & ? & \textbf{B} & \textbf{B} & U & \textbf{B}\\
Whisper Leak~\cite{whisperleak2025} & \textbf{B} & \textbf{B} & \textbf{B} & \textbf{P} & U & \textbf{B} & \textbf{B} & U & \textbf{P}\\
SpliceLeak~\cite{spliceleak2026} & \textbf{B} & \textbf{B} & \textbf{B} & -- & -- & \textbf{B} & \textbf{B} & U & \textbf{B}\\
Time Will Tell~\cite{timewilltell2024} & \textbf{B} & \textbf{B} & \textbf{B} & \textbf{P} & U & \textbf{B} & \textbf{B} & U & \textbf{B}\\
I Know What You Said~\cite{gao2025iknow} & \textbf{B} & \textbf{B} & \textbf{B} & -- & -- & \textbf{P} & -- & -- & \textbf{B}\\
Speculative decoding~\cite{wei2024speculation} & \textbf{B} & \textbf{B} & \textbf{P} & \textbf{P} & U & \textbf{B} & \textbf{B} & -- & \textbf{B}\\
Dangol workflow~\cite{dangol2026workflow} & \textbf{B} & \textbf{B} & \textbf{B} & \textbf{B} & \textbf{B} & \textbf{B} & \textbf{B} & U & \textbf{B}\\
Cloud traffic benchmark~\cite{pouryousef2026traffic} & \textbf{B} & \textbf{P} & \textbf{B} & U & U & \textbf{P} & \textbf{B} & \textbf{P} & U\\
Remote timing~\cite{carlini2024remote} & \textbf{B} & \textbf{B} & \textbf{P} & -- & -- & \textbf{B} & -- & U & \textbf{B}\\
TDXRay~\cite{hornetz2026tdxray} & \textbf{B} & \textbf{B} & \textbf{B} & -- & -- & \textbf{B} & -- & \textbf{B} & \textbf{B}\\
SafeKV~\cite{chu2025safekv} & \textbf{B} & \textbf{B} & \textbf{B} & -- & -- & \textbf{B} & \textbf{B} & U & \textbf{P}\\
AHB~\cite{zhu2026turing} & \textbf{B} & \textbf{B} & \textbf{B} & \textbf{P} & U & \textbf{B} & \textbf{B} & \textbf{B} & U\\
KVGov~\cite{addagada2026kvgov} & \textbf{B} & \textbf{B} & \textbf{B} & -- & -- & \textbf{B} & \textbf{B} & U & \textbf{B}\\
CachePrune~\cite{wu2026cacheprune} & \textbf{B} & \textbf{B} & \textbf{B} & -- & -- & \textbf{B} & \textbf{B} & \textbf{B} & \textbf{B}\\
OTRO~\cite{lee2026otro} & \textbf{B} & \textbf{B} & \textbf{B} & -- & -- & \textbf{B} & \textbf{B} & -- & \textbf{B}\\
SparLeak~\cite{chen2026sparleak} & \textbf{B} & \textbf{B} & \textbf{B} & \textbf{P} & U & \textbf{B} & \textbf{B} & \textbf{B} & U\\
Shape and Substance~\cite{hadad2026shape} & \textbf{B} & \textbf{B} & \textbf{P} & U & U & \textbf{B} & \textbf{P} & U & \textbf{P}\\
\midrule Bound & 20 & 17 & 17 & 4 & 2 & 17 & 16 & 4 & 11\\
Applicable & 20 & 20 & 20 & 11 & 11 & 20 & 17 & 17 & 19\\
\bottomrule\end{tabularx}
\par\smallskip\raggedright \textbf{B}: bound; \textbf{P}: partial; U: unreported in inspected scope; ?: unverified binding; --: inapplicable. Table~\ref{tab:audit-rules} defines R1--R9. Each cell is generated from the scoped records: B if any record is B, otherwise P, ?, U, A, then --. This availability summary retains a witness for every B; separate cells can refer to different configurations. NetEcho and Early Bird retain unresolved record-specific R5 bindings. Dangol's B in R5 is supported by stored capture; its live binding remains unresolved. S36 records the projection and its 13 corrections.
\end{table}


\soktakeaway{SoK Takeaway 2 (Inference is bound; reuse is not)}{Target and inference are bound almost everywhere (R1 20/20, R7 17/20). The bindings needed to carry a result elsewhere are not: protected fit 4/11 applicable families, matched adaptation 2/11 (2/10 for learned adaptive attacks only), utility at the protected operating point 3/17 (2/16 counting task utility only) and cost 11/19. CachePrune alone binds inference, utility and cost at one operating point, and it does not establish adaptation. These are scope sensitivities, not quality ranks.}
